\subsection{Classical vs.\ quantum: main comparison}
Table~\ref{tab:main-results} compares the classical and quantum branches
across all six scenarios. The quantum branch is less accurate in every
scenario (1.5--14.4$\times$ higher MAE) and substantially slower in every
scenario (53--89$\times$ wall-clock time), with no scenario favoring the
quantum branch on either axis.
% Table auto-generated from qcpinn_results/20260712_035925/ -- rerun
% tables/make_tables.py after any change to that CSV.
% Auto-generated by tables/make_tables.py -- do not edit by hand.
\begin{table}[t]
\centering
\caption{Classical vs.\ quantum PINN across six boundary-condition / problem-type
scenarios. MAE and wall-clock time (s) are means over seeds; ratio columns are
quantum relative to classical (higher = quantum worse).}
\label{tab:main-results}
\begin{tabular}{llrrrrrr}
\toprule
BC & Problem & MAE (cls.) & MAE (qtm.) & MAE ratio & Time (cls.) & Time (qtm.) & Time ratio \\
\midrule
Dirichlet & Forward & 0.000231 & 0.00034 & 1.5$\times$ & 31.9 & 2.44e+03 & 76.5$\times$ \\
Dirichlet & Inverse & 0.000287 & 0.0005 & 1.7$\times$ & 39.1 & 2.69e+03 & 69.0$\times$ \\
Neumann & Forward & 0.000166 & 0.00239 & 14.4$\times$ & 38.4 & 2.56e+03 & 66.6$\times$ \\
Neumann & Inverse & 0.000249 & 0.00118 & 4.7$\times$ & 35.6 & 2.82e+03 & 79.3$\times$ \\
Robin & Forward & 8.96e-05 & 0.000384 & 4.3$\times$ & 128 & 6.79e+03 & 52.9$\times$ \\
Robin & Inverse & 8.66e-05 & 0.00029 & 3.3$\times$ & 133 & 1.19e+04 & 89.3$\times$ \\
\bottomrule
\end{tabular}
\end{table}


\todo{Add 1-2 representative solution/error plots here, e.g.
\texttt{DE\_Dir\_FWD\_classical.png} vs \texttt{DE\_Dir\_FWD\_quantum.png}
from qcpinn\_results/20260712\_035925/ (already on the graphicspath in
main.tex), to show what the accuracy gap looks like spatially, not just in
the aggregate MAE number.}

\subsection{Ablation 1: training hyperparameters}
The main comparison leaves open whether the quantum branch's
underperformance reflects a suboptimal training recipe rather than a
structural limitation. We test two candidates on two representative
scenarios (Dirichlet-forward, Robin-inverse): a smaller weight
initialization scale (\texttt{small\_init}) intended to keep VQC rotation
angles away from regions with vanishing gradient, and a higher
learning-rate floor (\texttt{slow\_decay}) intended to prevent premature
convergence to a poor optimum. Table~\ref{tab:ablation1} shows neither
closes the gap: \texttt{slow\_decay} is the worst of the three quantum
configurations in both scenarios, and \texttt{small\_init} yields only a
3--16\% MAE improvement over the untouched baseline -- far short of the
3--14$\times$ gap to the classical branch.
% Auto-generated by tables/make_tables.py -- do not edit by hand.
\begin{table}[t]
\centering
\caption{Phase-1 ablation: quantum-branch initialization scale and LR-schedule
floor, vs.\ the classical reference and the untouched quantum baseline.
Neither intervention closes the accuracy or runtime gap.}
\label{tab:ablation1}
\begin{tabular}{llrr}
\toprule
Scenario & Config & MAE (mean $\pm$ std) & Time (s) \\
\midrule
Dirichlet Forward & classical & 0.000171 $\pm$ 3.7e-05 & 65.68 \\
Dirichlet Forward & baseline & 0.000522 $\pm$ 0.00013 & 5257 \\
Dirichlet Forward & small\_init & 0.000437 $\pm$ 6e-05 & 5097 \\
Dirichlet Forward & slow\_decay & 0.00102 $\pm$ 0.0008 & 5038 \\
Robin Inverse & classical & 7.91e-05 $\pm$ 2e-05 & 400.3 \\
Robin Inverse & baseline & 0.000626 $\pm$ 0.00033 & 1.492e+04 \\
Robin Inverse & small\_init & 0.000609 $\pm$ 0.00023 & 1.513e+04 \\
Robin Inverse & slow\_decay & 0.0011 $\pm$ 0.00032 & 1.476e+04 \\
\bottomrule
\end{tabular}
\end{table}


\subsection{Ablation 2: circuit structure}
Having ruled out these two hyperparameters, we next ask whether the circuit
itself is the bottleneck -- either under-expressive (too few qubits/layers)
or subject to a barren-plateau-like effect (gradients vanishing as the
circuit grows). We sweep qubit count and depth independently around the
baseline ($n{=}4, L{=}2$): \texttt{narrow} ($n{=}2$), \texttt{wide}
($n{=}6$), \texttt{deep} ($L{=}4$), holding initialization and
learning-rate schedule at the phase-1 baseline throughout. Alongside the
usual accuracy metrics we log the circuit-only gradient norm at the end of
training. Table~\ref{tab:ablation2} summarizes the result.
% Auto-generated by tables/make_tables.py -- do not edit by hand.
\begin{table}[t]
\centering
\caption{Phase-2 ablation: circuit width/depth sweep (narrow/baseline/wide/deep)
with final circuit-only gradient norm. Accuracy and gradient norm do not move
monotonically together, so neither a pure expressivity nor a pure
barren-plateau account fits cleanly -- see Discussion.}
\label{tab:ablation2}
\begin{tabular}{llrrr}
\toprule
Scenario & Config & MAE (mean $\pm$ std) & Final $\|\nabla_{\text{circuit}}\|$ & Time (s) \\
\midrule
Dirichlet Forward & classical & 0.000184 $\pm$ 3.5e-05 & 0 & 31.96 \\
Dirichlet Forward & baseline & 0.00052 $\pm$ 0.00013 & 1.56e-05 & 2655 \\
Dirichlet Forward & narrow & 0.00152 $\pm$ 0.00052 & 1.93e-05 & 968.1 \\
Dirichlet Forward & wide & 0.000446 $\pm$ 8.9e-05 & 9.45e-06 & 7398 \\
Dirichlet Forward & deep & 0.000663 $\pm$ 0.00023 & 2.21e-05 & 4431 \\
Robin Inverse & classical & 6.44e-05 $\pm$ 1.6e-05 & 0 & 151.3 \\
Robin Inverse & baseline & 0.000627 $\pm$ 0.00033 & 0.00653 & 8331 \\
Robin Inverse & narrow & 0.00177 $\pm$ 0.00092 & 0.0326 & 3173 \\
Robin Inverse & wide & 0.000413 $\pm$ 0.00013 & 0.00575 & 2.201e+04 \\
Robin Inverse & deep & 0.000508 $\pm$ 0.0001 & 0.00595 & 1.393e+04 \\
\bottomrule
\end{tabular}
\end{table}


\todo{Add the ablation2\_accuracy.png / ablation2\_qgrad.png figures from
qcpinn\_results/ablation2\_20260729\_225346/ here as a side-by-side, since
the table alone makes the non-monotonic pattern harder to see than the
plots do.}

The pattern is mixed rather than clean. \texttt{wide} achieves the best
accuracy among the quantum configurations in both scenarios, which would
suggest expressivity is a limiting factor worth scaling past -- but it also
has the \emph{smallest} final gradient norm of the four configurations,
which is not the direction a pure expressivity story predicts and is only
partially consistent with a barren-plateau account (which predicts smaller
gradients accompanying \emph{worse}, not better, accuracy). \texttt{narrow}
is the worst-performing configuration on accuracy, consistent with an
under-expressivity explanation at the low end, while \texttt{deep} is worse
than baseline on accuracy despite a comparable gradient norm. \todo{State
explicitly in the text -- and defend in Discussion -- that this rules out a
single clean explanation and motivates the more cautious framing adopted
there (data encoding / ansatz choice as an unresolved third factor).}
